\tzpanchor{0402}
\tzpentryheader{0402}{Particle-Specific Entanglement Formulas}

\begingroup\small
\begingroup\scriptsize
\setlength{\tabcolsep}{4pt}
\renewcommand{\arraystretch}{0.92}
\noindent\begin{tabularx}{\linewidth}{@{}p{0.24\linewidth}p{0.36\linewidth}X@{}}
\textbf{UUID} & \multicolumn{2}{l@{}}{\tzpcode{bc5153ed-f983-5593-bb15-dbf70af67672}}\\
\textbf{PiID} & \tzpcode{0572703657595} & \textbf{TZPi}\quad \tzpcode{3}\quad \textbf{PiPE}\quad \tzpcode{409}\\
\textbf{RTEID} & \textbf{Poloidal $\theta$}\quad \tzpcode{1561.0553 rad} & \textbf{Toroidal $\phi$}\quad \tzpcode{02.5721 rad}\\
\textbf{TZPID} & \tzpcode{ID0402} & \textbf{Node Dependencies}\quad \tzpidref{0401}\\
\textbf{Registry} & \tzpcode{registry\_id\_0402} & \textbf{Scale}\quad Quantum-Classical Transition\\
\textbf{LOC Classification} & QA641 manifolds and topology & QC174.17 mathematical physics\\
\textbf{arXiv Classification} & Primary: math-ph & Cross-list: gr-qc, math.DG\\
\textbf{Primary Support} & Canonical Definition & Support Relation\\
\textbf{Closure Support} & Closure Relation & \\
\end{tabularx}\par
\endgroup
\endgroup

\tzpsection{Canonical Statement}
ID 402 reifies the abstract Fneural into specific coupling Hamiltonians. This ensures the Wwitness can verify the "Depth of Handshake."

\tzpsection{Canonical Equation}
\[
\begin{adjustbox}{max width=\linewidth,center}
$\displaystyle
\Psi_{N}(\lambda_{0}) = \cdots
$
\end{adjustbox}
\]

\begin{minipage}[t]{0.49\linewidth}
\tzpsection{Canonical Symbol Key}
\begingroup
\small
\raggedright
\textbf{Source equation symbols.}\par
\begin{itemize}[leftmargin=1.35em,itemsep=1pt,topsep=1pt,parsep=0pt]
    \item $\Psi_{N}$: source equation symbol.
    \item $N$: normalization, count, or closure parameter in the entry relation.
    \item $\Gamma$: closed loop, path, or admissible trajectory in the entry domain.
    \item $\Lambda$: lemniscatic representative or canonical branch object.
    \item $\Psi$: wavefunction or state carrier for the entry equation.
    \item $\cdot s$: source equation symbol.
    \item $P$: source equation symbol.
    \item $i_{N}$: source equation symbol.
\end{itemize}
\endgroup
\end{minipage}\hfill
\begin{minipage}[t]{0.47\linewidth}
\tzpsection{Wolfram Form}
\begingroup
\scriptsize
\raggedright
\textbf{Source Wolfram model.}\par
\begin{Verbatim}[fontsize=\tiny,breaklines=true,breakanywhere=true]
(* TZPID ID0402 generated source Wolfram model *)
ClearAll[K0402, Cr0402, T, R, A, W, I, N];
eq0402$1 = HoldForm[PsiN[lambda0] == Ellipsis];

K0402 = HoldForm[{eq0402$1}];

N = "normalized TRAWIN closure object for Particle-Specific Entanglement Formulas";
I = "Particle-Specific Entanglement Formulas integration/comparison layer";
W = "Particle-Specific Entanglement Formulas wave or operator carrier";
A = "Particle-Specific Entanglement Formulas admissible action/amplitude condition";
R = "Particle-Specific Entanglement Formulas rotation/curl preservation layer";
T = "Particle-Specific Entanglement Formulas local source condition";

Cr0402[expr_] := HoldForm[N[I[W[A[R[T[expr]]]]]]];

Result0402 = Cr0402[K0402];
\end{Verbatim}
\endgroup
\end{minipage}

\par\vspace{0.45\baselineskip}
\noindent\begin{minipage}[t]{\linewidth}
\begingroup
\small
\raggedright
\textbf{TRAWIN Operator Key.}\par
\noindent\textbf{Time/localization operator:} $\mathsf{T}\!\left[\mathrm{local}\right]$ Particle-Specific Entanglement Formulas local source condition.\par
\noindent\textbf{Rotation/curl operator:} $\mathsf{R}\!\left[\nabla\times\right]$ Particle-Specific Entanglement Formulas rotation/curl preservation layer.\par
\noindent\textbf{Action/amplitude operator:} $\mathsf{A}\!\left[\mathrm{admissible}\right]$ Particle-Specific Entanglement Formulas admissible action/amplitude condition.\par
\noindent\textbf{Wave/d'Alembertian operator:} $\mathsf{W}\!\left[\Box\right]$ Particle-Specific Entanglement Formulas wave or operator carrier.\par
\noindent\textbf{Integration operator:} $\mathsf{I}\!\left[\int\right]$ Particle-Specific Entanglement Formulas integration/comparison layer.\par
\noindent\textbf{Normalization operator:} $\mathsf{N}\!\left[\mathrm{normalized}\right]$ normalized TRAWIN closure object for Particle-Specific Entanglement Formulas.\par
\endgroup
\end{minipage}

\clearpage
\noindent\textbf{[ID 0402] Particle-Specific Entanglement Formulas}\par
\vspace{0.35em}
\tzpsection{Derivation Layer}
\paragraph{Primary Statement.}
For any closed loop $\GammaLoop$ encircling the TZP in $M$, there exists a homotopy deforming $\GammaLoop$ into a lemniscate whose self-intersection occurs at $p$.

\paragraph{Primary Derivation.}
Let $\GammaLoop:S^1\to M$ be a closed loop surrounding $p$. Since $p$ is the limiting degeneracy point, paths enclosing $p$ may be continuously deformed while preserving their winding relation to $p$.

Introduce a homotopy:

\begin{equation}
H:S^1\times[0,1]\to M .
\end{equation}

At $s=0$, the homotopy recovers the original loop:

\begin{equation}
H(\theta,0)=\GammaLoop(\theta).
\end{equation}

At $s=1$, the loop is deformed into a figure-eight curve:

\begin{equation}
H(\theta,1)=\Lambda(\theta),
\end{equation}

where $\Lambda$ is a lemniscatic curve satisfying:

\begin{equation}
\Lambda(\theta_1)=\Lambda(\theta_2)=p,
\qquad
\theta_1\neq\theta_2 .
\end{equation}

Thus the self-intersection of the limiting loop is identified with the TZP.

\paragraph{Interpretation.}
The lemniscate is the canonical shape of TZP recurrence: two lobes, one crossing, one central point. This turns the zero point into a geometric crossing rather than a simple endpoint.

\par\vspace{0.35em}
\noindent\textbf{Derivable Consequences.}\quad
\begingroup
\footnotesize
\raggedright
The canonical equation anchors Particle-Specific Entanglement Formulas by connecting the source condition to the derivation layer and proof certificate.
\par\endgroup

\tzpsection{Equation Support}
\noindent\textbf{Canonical Support Relation}\par
\[
\begin{adjustbox}{max width=\linewidth,center}
$\displaystyle
\Psi_{N}(\lambda_{0}) = \cdots
$
\end{adjustbox}
\]
\noindent The support equation repeats the canonical relation so the derivation page remains self-contained.\par

\clearpage
\noindent\textbf{[ID 0402] Particle-Specific Entanglement Formulas}\par
\vspace{0.35em}
\tzpsection{Dictionary}
\begingroup\small
\tzpsection{Definition}
Particle-Specific Entanglement Formulas is a registry definition in the active TZPID arc.

\tzpsection{Contextual Meaning}
ID 402 reifies the abstract Fneural into specific coupling Hamiltonians. This ensures the Wwitness can verify the "Depth of Handshake."

\tzpsection{Formal Statement}
\[
\begin{adjustbox}{max width=\linewidth,center}
$\displaystyle
\Psi_{N}(\lambda_{0}) = \cdots
$
\end{adjustbox}
\]

\tzpsection{Technical Interpretation}
ID 402 reifies the abstract Fneural into specific coupling Hamiltonians. This ensures the Wwitness can verify the "Depth of Handshake."

\tzpsection{Equation Grounding}
The equation grounding is carried by the canonical equation and the source Wolfram model on the entry page.

\tzpsection{Interpretive Role}
The canonical equation anchors Particle-Specific Entanglement Formulas by connecting the source condition to the derivation layer and proof certificate.

\tzpsection{TZP Thesaurus}
Particle-Specific Entanglement Formulas; canonical operator; source equation; TRAWIN-normalized relation.
\endgroup

\tzpsection{Encyclopedia Mathematica}
\begingroup\footnotesize\sloppy
The visual plate below is reserved for the source-specific equation and progression graphic for [ID 0402]. It is included only as a companion to the canonical equation, not as a substitute for the formal text.
\par\endgroup
\begingroup\footnotesize\itshape
Visual plate pending source-specific approval for [ID 0402].
\par\endgroup

\clearpage
\begingroup\tzpsupportsetup
\noindent\textbf{\fontsize{10.8pt}{12.8pt}\selectfont [ID 0402] Constructive Logic and Proof Certificate}\par
\noindent\textbf{\fontsize{10pt}{12pt}\selectfont Particle-Specific Entanglement Formulas}\par
\vspace{0.45em}
\begingroup
\footnotesize
\setlength{\parskip}{0.22em}
\setlength{\parindent}{0pt}
\noindent\textbf{Curry--Howard--Lambek Interpretation}\par
Under the Curry--Howard--Lambek correspondence, this registry entry is read as a typed proof
object: the stated source condition supplies the proposition, the derivation supplies the
morphism, and the proof certificate records the machine handles needed to check the obligation
in the formal registry.

\vspace{0.45em}
\noindent\textbf{CHL Proof}\par
\textbf{Statement:} entanglementS = -sum\textbackslash{}\_k left[ n\textbackslash{}\_k ln n\textbackslash{}\_k + (1+n\textbackslash{}\_k)ln(1+n\textbackslash{}\_k) right]

\textbf{Logic Validation:}\\
\textbf{Proposition:} ``PARTICLE-SPECIFIC ENTANGLEMENT FORMULAS is valid as a tensorial/geometric statement when
the stated tensor fields are well formed on the specified domain.''\\
\textbf{Proof:} The source statement gives a metric, curvature, or tensor relation. The CHL reading
validates it by checking that the tensorial objects have compatible domains, indices,
and transformation behavior.

\textbf{VERDICT:} \ensuremath{\checkmark}\ \textbf{TRUE} --- tensorial well-formedness validation

\textbf{Formal Status:} \ensuremath{\checkmark}\ THEORETICAL -- source-backed CHL proof obligation

\textbf{Physical Status:} THEORETICAL -- requires model-specific empirical interpretation
\par\vspace{0.45em}
\noindent{\tiny\textbf{Isabelle/HOL.} \tzpplaincode{physics\_validity\_from\_model, external\_validity\_package, registry\_number\_matches, equation\_count\_matches}}\par
\vfill
\begingroup\footnotesize\itshape
Proof visual pending source-specific approval for [ID 0402].
\par\endgroup
\vfill
\endgroup
\endgroup
